\begin{ex}
Có 5 bạn học sinh đứng thành một hàng ngang: An, Bình, Cúc, Dũng, Hoa (tính từ trái sang phải). Bạn nào đứng ở chính giữa hàng? (An: 1, Bình: 2, Cúc: 3, Dũng: 4, Hoa: 5).
\shortans{3}
\loigiai{
Hàng có 5 bạn, bạn đứng ở chính giữa là bạn thứ ba: bạn Cúc (ghi số 3).
\textbf{Đáp số:} $3$.
}
\end{ex}

\begin{ex}
Tính từ trái sang phải, bạn đứng ngay bên trái bạn thứ ba là bạn đứng ở vị trí thứ mấy?
\shortans{2}
\loigiai{
Bạn đứng ngay bên trái bạn thứ ba là bạn ở vị trí thứ $2$.
\textbf{Đáp số:} $2$.
}
\end{ex}

\begin{ex}
Tính từ trái sang phải, bạn đứng ngay bên phải bạn thứ ba là bạn đứng ở vị trí thứ mấy?
\shortans{4}
\loigiai{
Bạn đứng ngay bên phải bạn thứ ba là bạn ở vị trí thứ $4$.
\textbf{Đáp số:} $4$.
}
\end{ex}

\begin{ex}
Một giá sách có 3 tầng: tầng trên, tầng giữa và tầng dưới. Tầng trên có 3 quyển sách, tầng giữa có 4 quyển, tầng dưới có 3 quyển. Hỏi cả giá sách có bao nhiêu quyển sách?
\shortans{10}
\loigiai{
Tổng số sách trên cả 3 tầng là: $3 + 4 + 3 = 10$ (quyển sách).
\textbf{Đáp số:} $10$.
}
\end{ex}

\begin{ex}
Trong một hàng dọc có 7 bạn học sinh. Bạn Nam đứng trước bạn Hoa và phía sau có 3 bạn. Hỏi có bao nhiêu bạn đứng trước bạn Nam?
\shortans{3}
\loigiai{
Có 3 bạn đứng sau Nam, cộng với Nam là 4 bạn. Số bạn đứng trước Nam là: $7 - 4 = 3$ (bạn).
\textbf{Đáp số:} $3$.
}
\end{ex}

\begin{ex}
Trên một chiếc bàn: bên trái có 4 quả táo, bên phải có 5 quả lê, ở giữa có 1 quả dưa hấu. Hỏi trên bàn có tất cả bao nhiêu quả?
\shortans{10}
\loigiai{
Tổng số quả trên bàn là: $4 + 5 + 1 = 10$ (quả).
\textbf{Đáp số:} $10$.
}
\end{ex}

\begin{ex}
Một cây cột đèn giao thông có 3 đèn màu: Đèn đỏ ở trên cùng, đèn vàng ở giữa, đèn xanh ở dưới cùng. Hỏi có bao nhiêu đèn màu trên cột đèn?
\shortans{3}
\loigiai{
Cột đèn có 3 đèn màu: đỏ, vàng, xanh (tất cả 3 đèn).
\textbf{Đáp số:} $3$.
}
\end{ex}

\begin{ex}
Bạn Hùng có hai bàn tay: bàn tay trái và bàn tay phải. Mỗi bàn tay có 5 ngón. Hỏi cả hai bàn tay có bao nhiêu ngón tay?
\shortans{10}
\loigiai{
Tổng số ngón tay của hai bàn tay là: $5 + 5 = 10$ (ngón).
\textbf{Đáp số:} $10$.
}
\end{ex}

\begin{ex}
Một đàn vịt bơi nối đuôi nhau gồm 6 con. Con vịt đầu đàn bơi trước con vịt cuối đàn bao nhiêu con vịt?
\shortans{5}
\loigiai{
Giữa con đầu tiên và con thứ sáu có 4 con, tính cả khoảng cách là: $6 - 1 = 5$ vị trí.
\textbf{Đáp số:} $5$.
}
\end{ex}

\begin{ex}
Trong một hàng có 8 bạn học sinh. Bạn đứng thứ 4 tính từ trái sang phải có bao nhiêu bạn đứng ở bên phải mình?
\shortans{4}
\loigiai{
Số bạn đứng ở bên phải là: $8 - 4 = 4$ (bạn).
\textbf{Đáp số:} $4$.
}
\end{ex}

\begin{ex}
Tính từ phải sang trái trong hàng gồm 9 bạn, bạn đứng ở vị trí thứ 6 là bạn đứng ở vị trí thứ mấy tính từ trái sang phải?
\shortans{4}
\loigiai{
Vị trí tính từ trái sang phải là: $9 - 6 + 1 = 4$.
\textbf{Đáp số:} $4$.
}
\end{ex}

\begin{ex}
Bên trái hộp có 3 chiếc bút chì, bên phải hộp có nhiều hơn bên trái 2 chiếc bút chì. Hỏi cả hai bên hộp có bao nhiêu chiếc bút chì?
\shortans{8}
\loigiai{
Bên phải có số bút là: $3 + 2 = 5$ (chiếc).\nCả hai bên có số bút là: $3 + 5 = 8$ (chiếc bút chì).
\textbf{Đáp số:} $8$.
}
\end{ex}

\begin{ex}
Một tòa nhà có 5 tầng. Bạn Nam ở tầng 4, bạn Bình ở tầng ngay dưới tầng của bạn Nam. Hỏi bạn Bình ở tầng mấy?
\shortans{3}
\loigiai{
Tầng ngay dưới tầng 4 là tầng $3$.
\textbf{Đáp số:} $3$.
}
\end{ex}

\begin{ex}
Một tòa nhà có 5 tầng. Bạn Mai ở tầng 2, bạn Hoa ở tầng ngay trên tầng của bạn Mai. Hỏi bạn Hoa ở tầng mấy?
\shortans{3}
\loigiai{
Tầng ngay trên tầng 2 là tầng $3$.
\textbf{Đáp số:} $3$.
}
\end{ex}

\begin{ex}
Trong một cuộc đua xe đạp, bạn Tuấn chạy ngay sau bạn thứ nhất và chạy ngay trước bạn thứ ba. Hỏi bạn Tuấn đang chạy ở vị trí thứ mấy?
\shortans{2}
\loigiai{
Chạy sau người thứ nhất và trước người thứ ba là chạy ở vị trí thứ $2$.
\textbf{Đáp số:} $2$.
}
\end{ex}

\begin{ex}
Có 4 chú chim đậu thành hàng ngang trên dây điện. Thêm 1 chú chim đậu vào bên trái và 2 chú chim đậu vào bên phải. Hỏi lúc này có bao nhiêu chú chim trên dây điện?
\shortans{7}
\loigiai{
Tổng số chim lúc này là: $4 + 1 + 2 = 7$ (chú chim).
\textbf{Đáp số:} $7$.
}
\end{ex}

\begin{ex}
Một cây xanh có cành phía trên có 5 quả xoài, cành phía dưới có 4 quả xoài. Hỏi cành phía trên nhiều hơn cành phía dưới bao nhiêu quả xoài?
\shortans{1}
\loigiai{
Cành trên nhiều hơn cành dưới là: $5 - 4 = 1$ (quả xoài).
\textbf{Đáp số:} $1$.
}
\end{ex}

\begin{ex}
Trong giờ thể dục, 10 bạn xếp hàng ngang. Thầy giáo gọi bạn đứng ở vị trí thứ 5 bước lên trước một bước. Hỏi còn lại bao nhiêu bạn đứng trong hàng?
\shortans{9}
\loigiai{
Số bạn còn lại trong hàng là: $10 - 1 = 9$ (bạn).
\textbf{Đáp số:} $9$.
}
\end{ex}

\begin{ex}
Một con rùa bò từ gốc cây sang bên phải 4 bước, sau đó quay lại bò sang bên trái 2 bước. Hỏi lúc này con rùa cách gốc cây bao nhiêu bước?
\shortans{2}
\loigiai{
Khoảng cách từ rùa đến gốc cây là: $4 - 2 = 2$ (bước).
\textbf{Đáp số:} $2$.
}
\end{ex}

\begin{ex}
Một chú ếch nhảy từ bờ ao sang bên phải 6 bước, rồi nhảy tiếp sang bên phải 3 bước nữa. Hỏi chú ếch đã nhảy cách bờ ao bao nhiêu bước?
\shortans{9}
\loigiai{
Tổng số bước chú ếch đã nhảy là: $6 + 3 = 9$ (bước).
\textbf{Đáp số:} $9$.
}
\end{ex}

\begin{ex}
Trong một bức tranh vẽ: mặt trời ở phía trên, ngôi nhà ở phía dưới và dòng sông ở ngoài cùng. Có bao nhiêu sự vật được kể tên?
\shortans{3}
\loigiai{
Có 3 sự vật là: mặt trời, ngôi nhà, dòng sông (tất cả 3 sự vật).
\textbf{Đáp số:} $3$.
}
\end{ex}

\begin{ex}
Bạn Hùng đứng đối diện với chiếc gương. Tay phải của Hùng cầm 1 quả táo. Hỏi trong gương, hình ảnh của bạn Hùng cầm quả táo ở phía tay nào? (Tay trái: ghi 1; Tay phải: ghi 2).
\shortans{1}
\loigiai{
Khi nhìn vào gương, hình ảnh bị đổi chiều trái - phải, nên tay phải của người sẽ thành tay trái trong gương (ghi số 1).
\textbf{Đáp số:} $1$.
}
\end{ex}

\begin{ex}
Có 3 hộp quà xếp chồng lên nhau: hộp đỏ ở dưới cùng, hộp xanh ở giữa, hộp vàng ở trên cùng. Hỏi hộp nào ở trên cùng? (Đỏ: 1, Xanh: 2, Vàng: 3).
\shortans{3}
\loigiai{
Hộp vàng ở trên cùng (ghi số 3).
\textbf{Đáp số:} $3$.
}
\end{ex}

\begin{ex}
Bàn học có 4 chân bàn: 2 chân phía trước và 2 chân phía sau. Hỏi chiếc bàn học có bao nhiêu chân bàn?
\shortans{4}
\loigiai{
Chiếc bàn học có tất cả: $2 + 2 = 4$ (chân bàn).
\textbf{Đáp số:} $4$.
}
\end{ex}

\begin{ex}
Một hàng cây có 5 cây trồng thẳng hàng cách đều nhau. Hỏi có bao nhiêu khoảng cách ở giữa các cây?
\shortans{4}
\loigiai{
Số khoảng cách giữa 5 cây trồng thẳng hàng là: $5 - 1 = 4$ (khoảng cách).
\textbf{Đáp số:} $4$.
}
\end{ex}
